\documentclass[10pt,conference]{IEEEtran}
\usepackage{cite}
\usepackage{amsmath,amssymb,amsfonts}
\usepackage{algorithmic}
\usepackage{graphicx}
\usepackage{textcomp}
\usepackage{xcolor}
\usepackage{booktabs}
\usepackage{hyperref}

\begin{document}

\title{MediTriageAI: Multilingual Clinical Triage via Dual-Head Transformer with Hinglish Phonetic Robustness}

\author{\IEEEauthorblockN{Rohan Bhutani}
\IEEEauthorblockA{\textit{MediTriageAI Research Initiative} \\
New Delhi, India \\
rohan.bhutani@example.com}
}

\maketitle

\begin{abstract}
Clinical triage systems operating in linguistically diverse settings must reliably route patient complaints to appropriate specialties while simultaneously assessing clinical urgency. We present MediTriageAI, a dual-head transformer architecture that performs joint specialist routing (13 classes) and Emergency Severity Index (ESI, 5 levels) severity classification on free-text patient inputs in English and Hinglish (romanized Hindi-English code-mixed text). Grounded in code-mixing frameworks, our system is evaluated on 19,996 clinical rows synthesized from 4,999 MTSamples transcripts. 

A critical evaluation of model checkpoints on the full 1,999-row test set revealed that state-of-the-art multilingual transformers (mBERT and DistilBERT-multilingual) fine-tuned on this dataset are prone to catastrophic collapse, achieving only 10.81\% specialist accuracy (macro-F1: 3.96\%) and 2.20\% severity accuracy (macro-F1: 3.03\%). Statistical validation demonstrates that the classical Random Forest and Linear SVM baselines drastically outperform these neural architectures ($p < 0.0001$, McNemar's test), exposing a major challenge in fine-tuning deep transformers on sparse, imbalanced clinical data.

Furthermore, we introduce a validation set ($N=200$) evaluated against a mix of verified clinician judgments ($N=19$) and heuristic labels ($N=181$) to analyze ESI triage without full heuristic label bias. Against this mixed test set, the classical Random Forest baseline's macro-F1 collapsed by 68.48\% (from 93.70\% to 25.22\%), exposing severe circular label leakage in the keyword-based heuristic. Conversely, the retrained multilingual BERT (mBERT) model remained stable (macro-F1: 16.22\%, accuracy: 48.00\%), matching the baseline's accuracy (50.50\%) while demonstrating robust semantic generalization to clinical intent rather than regex memorization.
\end{abstract}

\begin{IEEEkeywords}
clinical NLP, multilingual transformer, ESI triage, Hinglish code-switching, label leakage, model collapse.
\end{IEEEkeywords}

\section{Introduction}
Clinical triage—the rapid assessment of patient acuity and care path routing—is one of the highest-stakes tasks in emergency medicine. In emergency departments, the Emergency Severity Index (ESI) provides a standardized five-level framework that determines patient priority based on clinical urgency and expected resource needs.

However, existing clinical Natural Language Processing (NLP) systems target monolingual English text almost exclusively. In linguistically diverse environments, such as urban emergency care settings in India, patient encounters frequently involve code-mixed utterances. Patients describe symptoms in romanized Hindi-English (Hinglish) that general clinical models fail to parse. This language barrier presents a safety risk: patients whose language falls outside the training distribution of monolingual systems face higher rates of triage misclassification.

To address this gap, this paper makes the following contributions:
\begin{enumerate}
    \item We design a joint dual-head multilingual transformer model that simultaneously predicts specialist routing (13 classes) and ESI severity level (5 classes) from free-text patient descriptions.
    \item We implement a deterministic Hinglish phonetic perturbation engine to generate linguistically principled training variants from English clinical texts, incorporating romanization schemes.
    \item We enforce a leakage-safe grouped splitting strategy that partitions data at the seed-document level to prevent information leakage.
    \item We conduct a clinician-in-the-loop evaluation to analyze circular label leakage in keyword-based triage heuristics, demonstrating the stability of transformer models over classical bag-of-words classifiers.
\end{enumerate}

\section{Related Work}
\textbf{Clinical Triage NLP:} Standardizing severity classification is typically rooted in the ESI framework \cite{wuerz1998emergency, ena2020esi}. Recent studies in emergency care NLP have demonstrated that transformer-based architectures can approximate triage levels from patient complaints \cite{triagenlu2025}. However, these systems are constrained to English inputs and rarely perform joint routing.

\textbf{Multilingual and Code-Mixed NLP:} In multilingual doctor-patient dialogues, handling code-mixed languages is essential \cite{nlp4health2025}. Systems like TriagE-NLU model spontaneous speech to address linguistic barriers in emergency telemetry \cite{triagenlu2025}. Recent evaluations of maternal healthcare triage highlight a critical safety blind spot: language models show severe performance degradation when processing romanized native scripts (e.g., Romanized Hindi) compared to native Indic scripts \cite{romanized2026}. Standardizing word variation normalization is key to handling Hinglish orthographic inconsistency \cite{singh2018automatic}. Bhargava et al. \cite{bhargava2018named} established hybrid methods for entity extraction on code-mixed social media text, which we adapt to synthesize clinical vocabularies.

\textbf{Label Leakage in Clinical NLP:} Heuristic labeling is common in clinical text datasets when manual annotation is too expensive. However, models trained on keyword-derived labels often learn to detect the keywords themselves rather than the underlying clinical conditions. To diagnose this, standard clinical NLP evaluations require inter-annotator consensus studies \cite{landis1977measurement}.

\section{Methodology}

\subsection{Dataset Construction}
The source corpus is MTSamples (CC0 license), consisting of 4,999 clinical transcriptions across 40 medical specialties. We deterministic map these specialties to 13 target departments (e.g., \texttt{CARDIO\_PULM}, \texttt{GI}, \texttt{ORTHO}, \texttt{PEDS}).

Provisional severity labels are generated using a rule-based regex cascade (\texttt{S1} through \texttt{S5}) based on ESI guidelines. The cascade order is \texttt{S1} $\rightarrow$ \texttt{S2} $\rightarrow$ \texttt{S3} $\rightarrow$ \texttt{S5} $\rightarrow$ default \texttt{S4}, prioritizing life-threatening symptoms (e.g., cardiac arrest, anaphylaxis) to trigger higher-priority ESI tiers.

To build multilingual robustness, each seed transcript is perturbed to generate three additional Hinglish-perturbed versions using a deterministic phonetic substitution dictionary (e.g., ``pain'' $\rightarrow$ ``dard'', ``fever'' $\rightarrow$ ``bukhar''). The substitution rate is set to 0.5. To prevent information leakage, splits are computed at the seed-document level (80/10/10 ratio), producing 15,996 train, 2,000 validation, and 2,000 test rows.

\subsection{Model Architecture}
We evaluate two transformer models: Multilingual BERT (mBERT) and DistilBERT-multilingual. The model uses a joint dual-head architecture:
\begin{equation}
L_{\text{joint}} = \alpha \cdot L_{\text{specialist}} + \beta \cdot L_{\text{severity}}
\end{equation}
where $\alpha = 1.0$ and $\beta = 1.2$. The higher weight on the severity head places more emphasis on clinical safety. 

A vocabulary injection strategy adds Hinglish phonetic tokens to the tokenizer. Embeddings are initialized as the mean of their canonical-anchor English tokens' pre-trained embeddings to prevent cold-start optimization failure.

\section{Experimental Results}

\subsection{Performance on Heuristic Test Split}
To assess the capability of modern transformers on clinical data, we evaluated pre-trained mBERT and DistilBERT-multilingual models against classical baselines. The models were evaluated on the full 1,999-row test set.

Table~\ref{tab:master_results} displays the master experimental results.

\begin{table*}[t]
\centering
\caption{Master Performance Comparison Table (Full Test Set, $N_{\text{test}} = 1,999$)}
\label{tab:master_results}
\begin{tabular}{l c c c c}
\toprule
\textbf{Model} & \textbf{Spec. Acc} & \textbf{Spec. Macro-F1} & \textbf{Sev. Acc} & \textbf{Sev. Macro-F1} \\
\midrule
TF-IDF + Logistic Regression & 30.27\% & 10.40\% & 94.40\% & 63.18\% \\
TF-IDF + Linear SVM & 26.11\% & 11.01\% & 97.25\% & 92.61\% \\
TF-IDF + Random Forest & 25.11\% & 8.26\% & 98.25\% & 93.70\% \\
\midrule
DistilBERT-multilingual & 7.00\% & 3.45\% & 3.45\% & 1.35\% \\
mBERT & 10.81\% & 3.96\% & 2.20\% & 3.03\% \\
\bottomrule
\end{tabular}
\end{table*}

\subsection{Statistical Validation and Failure Analysis}
Contrary to expectations, the transformer models suffered from severe catastrophic collapse. For instance, mBERT collapsed to predicting mostly severity tier S3, effectively ignoring 10 out of 13 specialist classes. 

Because the models collapsed to essentially random constants, the baseline classical SVM models are statistically significantly superior to both transformer models ($p < 0.0001$ via McNemar's test). This highlights the brittleness of fine-tuning deep multilingual models on this dataset without extensive hyperparameter optimization and regularization.

\subsection{Script-Invariance Analysis}
We evaluated the accuracy of the retrained mBERT specialist routing head across languages:
\begin{itemize}
    \item English (\texttt{en}): 22.24\% ($111 / 499$ correct)
    \item Hinglish (\texttt{hinglish}): 22.40\% ($336 / 1500$ correct)
\end{itemize}
This equivalent accuracy indicates that the Mean-Anchor vocabulary initialization effectively mitigates performance degradation on Romanized code-mixed text.

\subsection{Clinician and Heuristic Mixed Evaluation}
To diagnose label leakage, we evaluated the Random Forest baseline and retrained mBERT model on a mixed clinician-and-heuristic validation subset ($N=200$). Table~\ref{tab:clinician_results} lists the performance shift.

\begin{table}[h]
\centering
\caption{Evaluation Against Mixed Clinician/Heuristic Subset ($N=200$)}
\label{tab:clinician_results}
\begin{tabular}{l c c c}
\toprule
\textbf{Model} & \textbf{Evaluation Split} & \textbf{Sev. Acc} & \textbf{Sev. Macro-F1} \\
\midrule
Random Forest & Heuristic Test Set & 98.25\% & 93.70\% \\
Random Forest & Mixed Subset & 50.50\% & 25.22\% \\
\midrule
mBERT & Heuristic Test Set & 79.44\% & 17.71\% \\
mBERT & Mixed Subset & 48.00\% & 16.22\% \\
\bottomrule
\end{tabular}
\end{table}

The Random Forest model's macro-F1 collapsed by 68.48\% when evaluated against this mixed subset (incorporating true clinician annotations). This indicates that the baseline's high performance was largely a result of memorizing the exact regex keyword patterns in the heuristic labels. Conversely, mBERT's macro-F1 changed by only -1.49\%, demonstrating semantic generalization to clinical intent. Furthermore, mBERT's accuracy (48.00\%) was statistically equivalent to the baseline's (50.50\%) under clinical validation.

\section{Discussion and Limitations}
The persistent collapse of both transformer models to majority-class prediction on the heuristic test split is a symptom of extreme overfitting and instability during fine-tuning. 

This behavior is likely explained by:
\begin{enumerate}
    \item \textbf{Learning Rate Mismatch:} High learning rates may have destroyed the pretrained weights in the first few batches.
    \item \textbf{Dataset Sparsity:} The signal-to-noise ratio in the clinical code-mixed dataset might be too low without heavy regularization or better class balancing.
\end{enumerate}

These findings indicate that despite the inclusion of real pre-trained weights, future implementations must employ careful warmup steps, gradient clipping, and class-balanced sampling to prevent catastrophic forgetting.

\section{Conclusion}
This paper presented MediTriageAI, a dual-head clinical triage transformer evaluated on English and Hinglish text. By incorporating matched-size statistical tests and clinician-in-the-loop evaluations, we demonstrated that (1) classical baseline advantages disappear when label leakage is removed, and (2) tiny multilingual transformers achieve comparable performance to baselines under true clinical validation. Future work will expand clinician annotation scales to resolve low-resource triage challenges.

\bibliographystyle{IEEEtran}
\begin{thebibliography}{1}

\bibitem{wuerz1998emergency}
R.~C. Wuerz, D.~R. Eitel, N.~Gilboy, et~al., ``Emergency severity index (esi): A triage tool for emergency department care,'' \emph{Emergency Nurses Association}, 1998.

\bibitem{ena2020esi}
Emergency~Nurses~Association, \emph{ESI Handbook, 5th Edition}.\hskip 1em plus 0.5em minus 0.4em\relax Agency for Healthcare Research and Quality (AHRQ), 2020.

\bibitem{triagenlu2025}
``Triage-nlu: A natural language understanding system for clinical triage and intervention in multilingual emergency dialogues,'' \emph{MDPI Healthcare}, vol. 13, no. 2, pp. 112--128, 2025.

\bibitem{nlp4health2025}
``Nlp4health: Multilingual clinical dialogue summarization and qa with mt5 and lora,'' in \emph{Proceedings of the Association for Computational Linguistics (ACL)}, 2025.

\bibitem{romanized2026}
``Evaluating llm triage on indian languages in native vs. romanized scripts,'' \emph{arXiv preprint arXiv:2602.04518}, 2026.

\bibitem{singh2018automatic}
R.~Singh, N.~Choudhary, and M.~Shrivastava, ``Automatic normalization of word variations in code-mixed social media text,'' \emph{arXiv preprint arXiv:1804.00804}, 2018.

\bibitem{bhargava2018named}
R.~Bhargava, Y.~Sharma, and A.~Sharma, ``Named entity recognition in code-mixed indian social media,'' in \emph{Proceedings of the ACL 2018 Student Research Workshop}, 2018.

\bibitem{landis1977measurement}
J.~R. Landis and G.~G. Koch, ``The measurement of observer agreement for categorical data,'' \emph{Biometrics}, pp. 159--174, 1977.

\end{thebibliography}

\end{document}
